% ===================================================================
%  TAVOLA N. 7 — Il paese d'inverno. --- La fattoria in primavera.
%  Traduction 2026 d'apres texto/io, controlee sur texto/fr et
%  texto/en. Consigne : voir texto/it/10-tavola-01.tex.
%
%  LE MOT « CONTADINO » DIT QUI A NOMME LA CAMPAGNE, ET CE N'EST PAS
%  LA CAMPAGNE.
%
%  Le paysan de l'alinea 23 de la premiere scene et celui de l'alinea
%  8 de la seconde s'appellent en italien contadino, du
%  « contado », qui vient de « comitatus » : LE TERRITOIRE SOUMIS A
%  UNE VILLE. Le contado de Florence, celui de Sienne, celui de
%  Bologne : ce sont des circonscriptions administratives, definies
%  depuis la ville et pour la ville.
%
%  Les autres langues nomment cet homme par ce qu'il fait ou par ou il
%  vit :
%     allemand  « Bauer »    celui qui batit, qui habite ;
%     francais  « paysan »   l'homme du pays ;
%     anglais   « farmer »   celui qui tient un bail ;
%     estonien  « talupoeg » le fils de la ferme.
%  L'ITALIEN SEUL LE NOMME PAR LA JURIDICTION D'UNE VILLE. Le mot
%  suppose qu'on regarde depuis les murs, et il n'a jamais cesse de le
%  supposer.
%
%  ET LA PLAISANTERIE DU CORDONNIER DE L'ALINEA 6 CHANGE DE CHARNIERE
%  UNE TROISIEME FOIS. Le vieux Giacomo dit : en ville j'etais l'un, ici
%  je suis l'autre.
%     francais  cordonnier / savetier   : deux METIERS, l'un fait,
%               l'autre repare ; et le premier est nomme par le cuir
%               de Cordoue ;
%     allemand  Schuhmacher / Schuster  : deux FORMES DU MEME MOT, la
%               longue et la courte ;
%     italien   calzolaio / ciabattino  : deux OBJETS — celui de la
%               chaussure et celui de la savate.
%  Trois langues, trois charnieres, et la meme plaisanterie tient dans
%  les trois. Ce n'est pas le mot qui la porte, c'est L'ECART entre
%  deux mots, quel qu'il soit.
%
%  ENFIN LE FATTORE DE L'ALINEA 8 EST « CELUI QUI FAIT », du latin
%  « factor » : c'est le regisseur, celui qui dirige le domaine pour
%  le compte d'un autre. A cote de lui l'italien a le mezzadro,
%  le metayer, mot a mot CELUI DE LA MOITIE — il gardait la moitie de
%  la recolte et rendait l'autre. La mezzadria a tenu de 1200 a 1964 ;
%  elle est la seule institution rurale d'Europe qui ait donne son nom
%  a l'homme qui la subissait, et non a la terre ni au contrat.
% ===================================================================

%%K t07-noto-1 noto
\VUnotes{143.2mm}{%
(*) Per i particolari di un pasto vedi le tavole 6 e 13.%
}

%%K t07-apar-1 apar
\VUsaut{4.0mm}
\VUcentre{13.2pt}{90}{SECONDA SERIE}
\VUsaut{3.0mm}
\VUfilet{16mm}
\VUsaut{5.6mm}
\VUtitre{15.0pt}{60}{TAVOLA N. 7}
\VUsaut{3.0mm}

%%K t07-tit-1 sub
\VUcentre{12.0pt}{0}{\textit{Prima scena.}}
\VUsaut{3.0mm}

%%K t07-tit-2 sub
\VUpk{10.2pt}{40}{Il paese d'inverno.}
\VUsaut{4.0mm}

%%K t07-c1-01-1 p
1. --- Durante le vacanze di Capodanno sono andato con mio padre a Villanova, un
paesino dove lui aveva delle cose da sbrigare.

%%K t07-c1-02-1 p
2. --- Abbiamo preso la \VUgras{diligenza} \textsuperscript{(11)} che fa servizio fra
la città e il paese; ma faceva molto freddo, e il viaggio su una vettura così scomoda
non è stato per niente piacevole.

%%K t07-c1-03-1 p
3. --- Siamo stati davvero contenti quando abbiamo visto il \VUgras{cocchiere}
fermare la vettura sulla piazza, davanti alla \VUgras{locanda} \textsuperscript{(2)}
dell'\textit{Orso Bianco}. L'\VUgras{oste} \textsuperscript{(4)} ci aspettava sulla
soglia, fumando tranquillamente la sua \VUgras{pipa}.

%%K t07-c1-03-2 p
Ci ha fatto entrare, e non me lo sono fatto dire due volte prima di sistemarmi davanti
al fuoco di sarmenti che scoppiettava nel camino. Dopo di che del \VUgras{vento del
nord} pungente non mi importava più niente, quel vento che faceva cigolare la vecchia
\VUgras{insegna} \textsuperscript{(3)} e portava via in volute nere il \VUgras{fumo}
\textsuperscript{(1)} che usciva dal camino.

%%K t07-c1-04-1 p
4. --- Era ora di pranzo. Senza aspettare oltre abbiamo fatto onore al pasto semplice
ma saporito che ci hanno servito (*).

%%K t07-tit-3 sub
\VUpk{10.2pt}{40}{Qualche visita.}
\VUsaut{3.0mm}

%%K t07-c1-05-1 p
5. --- Dopo pranzo sono andato con mio padre, che fa l'agente di assicurazioni, dai
suoi clienti. Prima siamo andati dal \VUgras{maniscalco} \textsuperscript{(5)}, che
abita accanto alla locanda. Stava ferrando un cavallo. Ho ammirato la forza del suo
garzone, che teneva lo \VUgras{zoccolo} del cavallo saldo contro il suo grembiule di
cuoio mentre il maniscalco fissava il \VUgras{ferro} \textsuperscript{(6)} a colpi di
martello.

%%K t07-c1-06-1 p
6. --- Poi siamo andati dal \VUgras{calzolaio} \textsuperscript{(7)} del paese, il
vecchio Giacomo. Anche se aveva per insegna un magnifico \VUgras{stivale}, si
accontentava di risuolare un vecchio paio di scarpe consumate.

%%K t07-c1-06-2 p
Il vecchio ci ha detto ridendo: «In città facevo il ,calzolaio‘ e non avevo clienti.
Qui faccio il ,ciabattino‘ e di lavoro ne ho quanto voglio.»

%%K t07-c1-07-1 p
7. --- Ci ha parlato del suo vicino, il \VUgras{barbiere} \textsuperscript{(8)}, di cui
tutti i clienti sono scontenti. I suoi \VUgras{rasoi} sono sempre intaccati e le sue
\VUgras{forbici} non tagliano mai bene. Ma siccome è l'unico barbiere del paese, alla
gente non resta che farsi radere e tagliare i capelli da lui. Per giunta è avaro e
antipatico.

%%K t07-c1-08-1 p
8. --- Per fortuna gli altri \VUgras{vicini} non sono come lui. Il
\VUgras{carradore} \textsuperscript{(9)}, per esempio, è un brav'uomo, lavoratore e
servizievole. È un piacere dargli un carro da riparare. Il suo lavoro è sempre finito
bene.

%%K t07-c1-09-1 p
9. --- Nemmeno del \VUgras{sellaio} \textsuperscript{(10)}, che ogni tanto aiuta a
cucire i \VUgras{finimenti} quando il lavoro si accumula, il vecchio Giacomo aveva da
lamentarsi.

%%K t07-c1-09-2 p
Mio padre gli ha chiesto della sua famiglia. «Stanno tutti bene. Mia nipote è a
\VUgras{scuola} \textsuperscript{(12)}. La sua \VUgras{maestra} \textsuperscript{(13)}
è molto contenta di lei; è una brava \VUgras{scolara} \textsuperscript{(14)}. Non è
come quel discolo di mio figlio, che pensa solo a giocare. Il \VUgras{maestro}
\textsuperscript{(15)} non riesce a fargli imparare niente.»

%%K t07-tit-4 sub
\VUsaut{3.0mm}
\VUpk{10.2pt}{40}{Chiacchiere di paese.}
\VUsaut{3.0mm}

%%K t07-c1-10-1 p
10. --- Poi il vecchio ci ha raccontato le novità del paese. Si doveva costruire una
scuola nuova, perché quella del \VUgras{municipio} era diventata troppo piccola.
Sarebbe stato facile: bastava comprare una parte del giardino del \VUgras{mugnaio}.
Al pover'uomo avrebbe fatto molto comodo. Da quando era arrivato il freddo, le
\VUgras{macine} del \VUgras{mulino} \textsuperscript{(16)} non potevano più macinare il
grano, perché la \VUgras{ruota} \textsuperscript{(17)} era presa nel
\VUgras{ghiaccio} \textsuperscript{(25)} del \VUgras{ruscello}
\textsuperscript{(23)}.

%%K t07-c1-11-1 p
11. --- «A proposito di costruzioni: ha visto la nostra nuova \VUgras{chiesa}
\textsuperscript{(18)}? Sotto il suo mantello di neve è molto pittoresca. Le
\VUgras{cornacchie} \textsuperscript{(19)}, che non riconoscono più il loro vecchio
campanile, stanno tristi sul grande albero del \VUgras{cimitero}
\textsuperscript{(20)}.»

%%K t07-c1-12-1 p
12. --- Il cimitero ha fatto venire in mente al calzolaio le persone che mio padre
aveva conosciuto e che erano morte nell'ultimo anno. «Quante \VUgras{tombe}
\textsuperscript{(21)}, caro signore, si sono aggiunte alle altre sotto il
\VUgras{cipresso} \textsuperscript{(22)}! Il nostro vecchio notaio ce l'ha portato via
la morte. Al suo \VUgras{funerale} c'era tutto il paese.»

%%K t07-c1-13-1 p
13. --- «Ecco lì il suo successore, che pattina sul ruscello. È passato poco fa perché
gli aggiustassi la cinghia del suo \VUgras{pattino} \textsuperscript{(26)}, che si era
rotta.»

%%K t07-c1-14-1 p
14. --- «E là, sulla riva del ruscello, c'è la nuova \VUgras{guardia campestre}
\textsuperscript{(27)}. È un ex carabiniere ed è il terrore dei monelli del paese.
Scappano appena vedono le sue \VUgras{ghette} \textsuperscript{(29)} e il suo
\VUgras{berretto d'ordinanza} \textsuperscript{(28)}.»

%%K t07-c1-15-1 p
15. --- «Ah, ecco il vecchio Giuseppe che porta un \VUgras{carico di fascine}
\textsuperscript{(30)} per il fornaio. --- Proprio così, ha detto mio padre, riconosco
la sua pariglia di \VUgras{muli} \textsuperscript{(31)}. Devo parlargli, e mi viene
proprio comodo. Arrivederci, Giacomo.»

%%K t07-c1-16-1 p
16. --- Proprio mentre uscivamo, i bambini del paese uscivano da scuola e si buttavano
sullo \VUgras{scivolo di ghiaccio} \textsuperscript{(33)}, con il rischio di cadere e
di rompersi qualcosa nel \VUgras{ruzzolone} \textsuperscript{(34)}. Uno di loro è
andato a prendere la sua \VUgras{slitta} \textsuperscript{(32)} dal carradore, ci ha
fatto sedere un amico ed è partito di corsa.

%%K t07-c1-17-1 p
17. --- Altri si sono precipitati verso un \VUgras{mucchio di neve}
\textsuperscript{(37)} vicino al \VUgras{ponte} \textsuperscript{(40)} e hanno preso a
bersaglio il \VUgras{pupazzo di neve} \textsuperscript{(39)} che avevano fatto il
giorno prima e a cui avevano messo un cappello, una pipa e una cartella a tracolla.

%%K t07-c1-18-1 p
18. --- Del vento gelido e del freddo non si curavano affatto. La maggior parte non
aveva nemmeno una \VUgras{sciarpa} \textsuperscript{(35)}, e per riscaldarsi dopo la
battaglia a \VUgras{palle di neve} \textsuperscript{(38)} bastava loro una manciata di
\VUgras{castagne} \textsuperscript{(42)} bollenti comprate dalla vecchia Giacomina, la
\VUgras{venditrice}, che trema dal freddo sotto il suo \VUgras{scialle}
\textsuperscript{(43)} troppo leggero.

%%K t07-tit-5 sub
\VUsaut{3.0mm}
\VUcentre{12.0pt}{0}{\textit{Seconda scena.}}
\VUsaut{3.0mm}

%%K t07-tit-6 sub
\VUpk{10.2pt}{40}{La fattoria in primavera.}
\VUsaut{4.0mm}

%%K t07-c2-01-1 p
1. --- L'inverno ha i suoi piaceri, ma il ritorno della \VUgras{primavera} lo salutiamo
lo stesso con grande gioia.

%%K t07-c2-01-2 p
Che bel quadro fa un'aia la sera, nel mese di maggio!

%%K t07-c2-02-1 p
2. --- All'orizzonte il \VUgras{sole} \textsuperscript{(1)} tramonta lentamente e
inonda la \VUgras{campagna} \textsuperscript{(3)} dei suoi \VUgras{raggi}
\textsuperscript{(2)} purpurei. L'\VUgras{aratore} \textsuperscript{(6)} operoso si
affretta ad arare il suo \VUgras{campo} \textsuperscript{(4)}.

%%K t07-c2-02-2 p
Con il suo \VUgras{aratro} \textsuperscript{(5)} tirato da un robusto \VUgras{cavallo}
\textsuperscript{(7)}, traccia il \VUgras{solco} \textsuperscript{(8)} in cui il
\VUgras{seminatore} \textsuperscript{(9)} spargerà a piene mani la \VUgras{semente} da
cui verranno le spighe dorate.

%%K t07-c2-03-1 p
3. --- I \VUgras{falciatori} \textsuperscript{(11)} con le loro falci hanno tagliato
l'erba del prato, i fienaioli hanno asciugato il \VUgras{fieno}
\textsuperscript{(10)} rivoltandolo con i \VUgras{forconi} \textsuperscript{(12)} e i
rastrelli, ed eccoli che tornano a casa.

%%K t07-c2-03-2 p
Alcuni camminano davanti al pesante carro tirato dai buoi, con gli attrezzi sulle
spalle. Altri, più pigri, si sono sistemati comodi sul fieno profumato.

%%K t07-c2-04-1 p
4. --- Passando salutano il \VUgras{giardiniere} \textsuperscript{(16)}, occupato
nel \VUgras{frutteto} \textsuperscript{(13)} a potare gli \VUgras{alberi da frutto} e a
innestare i \VUgras{peri} \textsuperscript{(14)}. I \VUgras{susini}
\textsuperscript{(15)} sono in fiore, e nell'\VUgras{orto} \textsuperscript{(17)} ben
concimato gli ortaggi, i cavoli e le lattughe vengono già su bene.

%%K t07-c2-04-2 p
Sul muretto un bel \VUgras{pavone} \textsuperscript{(18)} fa la ruota per mostrare le
sue magnifiche penne.

%%K t07-tit-7 sub
\VUpk{10.2pt}{40}{L'aia.}
\VUsaut{3.0mm}

%%K t07-c2-05-1 p
5. --- Le porte della \VUgras{stalla} \textsuperscript{(19)} sono aperte, e si vedono la
rastrelliera e la \VUgras{lettiera} \textsuperscript{(23)} della \VUgras{cavalla}
\textsuperscript{(21)} che il \VUgras{garzone di stalla} \textsuperscript{(20)} ha
appena staccato. Il \VUgras{puledro} \textsuperscript{(22)}, così buffo sulle sue lunghe
gambe, segue la madre facendo salti.

%%K t07-c2-06-1 p
6. --- Vicino al \VUgras{pozzo} \textsuperscript{(29)} le \VUgras{mucche}
\textsuperscript{(25)} vanno a bere al \VUgras{trogolo} \textsuperscript{(26)} di legno
che serve loro da abbeveratoio. Il \VUgras{vitello} \textsuperscript{(24)} ne
approfitta per poppare, e il \VUgras{bue} \textsuperscript{(27)}, che sente il buon
odore del foraggio, muggisce contento e alza fiero la testa con le sue potenti
\VUgras{corna} \textsuperscript{(28)}.

%%K t07-c2-07-1 p
7. --- Sono tutti al lavoro. La \VUgras{serva} \textsuperscript{(32)} tiene la
\VUgras{fune} \textsuperscript{(31)} del pozzo. Tira su l'acqua in un
\VUgras{secchio} \textsuperscript{(30)} per dare da bere agli animali. Vicino al
\VUgras{fienile} \textsuperscript{(33)} un uomo raccoglie con il rastrello la paglia
sparsa, e davanti alla porta della \VUgras{casa colonica} \textsuperscript{(34)} una
vecchia \VUgras{filatrice} \textsuperscript{(35)}, con il suo filatoio e la sua
\VUgras{rocca}, fila il \VUgras{lino} o la \VUgras{canapa} come una volta.

%%K t07-tit-8 sub
\VUsaut{3.0mm}
\VUpk{10.2pt}{40}{Il fattore.}
\VUsaut{3.0mm}

%%K t07-c2-08-1 p
8. --- Il \VUgras{fattore} \textsuperscript{(36)} dirige tutto questo lavoro e dà gli
ordini ai suoi uomini. Dice loro di mettere al riparo l'\VUgras{erpice}
\textsuperscript{(40)} e il \VUgras{piccone} \textsuperscript{(39)} rimasti vicino alla
\VUgras{cuccia} \textsuperscript{(38)} del \VUgras{cane} \textsuperscript{(37)}.

%%K t07-c2-09-1 p
9. --- Le sue grida fanno alzare in volo un \VUgras{piccione} \textsuperscript{(41)},
che scappa a tutta velocità verso la \VUgras{colombaia} \textsuperscript{(42)}, e le
\VUgras{galline} \textsuperscript{(43)}, che si affrettano a rientrare nel
\VUgras{pollaio} \textsuperscript{(44)}.

%%K t07-c2-10-1 p
10. --- Davanti all'\VUgras{ovile} \textsuperscript{(48)}, il vecchio \VUgras{pastore}
\textsuperscript{(45)} riporta il suo \VUgras{gregge} \textsuperscript{(49)}. Con il
suo \VUgras{vincastro} \textsuperscript{(46)} in mano, che non lo lascia mai più del
suo \VUgras{camiciotto} \textsuperscript{(47)} sbiadito, porta all'ovile
\VUgras{montoni} \textsuperscript{(50)}, \VUgras{pecore} \textsuperscript{(53)},
\VUgras{agnelli} \textsuperscript{(54)}, \VUgras{capre} \textsuperscript{(51)} e
\VUgras{capretti} \textsuperscript{(52)}. Il suo fedele \VUgras{cane}
\textsuperscript{(55)} Argo viene per ultimo e sprona con i suoi latrati chi resta
indietro.

%%K t07-c2-11-1 p
11. --- Tutto questo rumore e questo andirivieni non turbano la calma della buona
\VUgras{mucca da latte} \textsuperscript{(56)}, che fa suonare il suo
\VUgras{campanaccio} \textsuperscript{(57)} mentre la mungono davanti alla porta della
\VUgras{stalla delle mucche} \textsuperscript{(58)}.

%%K t07-c2-12-1 p
12. --- Sembra capire che deve dare il suo latte perché la \VUgras{massaia}
\textsuperscript{(60)} possa fare il \VUgras{burro} nella \VUgras{zangola}
\textsuperscript{(61)}, o gli ottimi \VUgras{formaggi} \textsuperscript{(64)} che
sgocciolano sull'asse posata sul \VUgras{mastello} \textsuperscript{(63)}.

%%K t07-c2-13-1 p
13. --- Tutta la \VUgras{casa del latte} \textsuperscript{(59)} è tenuta pulitissima
dal \VUgras{bifolco} \textsuperscript{(65)}, che ha appena lavato i \VUgras{bidoni del
latte} \textsuperscript{(62)} nel grande secchio che porta.

%%K t07-tit-9 sub
\VUsaut{3.0mm}
\VUpk{10.2pt}{40}{Il cortile del pollame.}
\VUsaut{3.0mm}

%%K t07-c2-14-1 p
14. --- Con i suoi pesanti zoccoli cammina in modo piuttosto goffo, e fa scappare il
\VUgras{tacchino} \textsuperscript{(68)}, le \VUgras{anatre} \textsuperscript{(66)}, i
\VUgras{maschi delle anatre} \textsuperscript{(67)} e le \VUgras{oche}
\textsuperscript{(69)}, che spalancano i loro \VUgras{becchi} \textsuperscript{(70)} e
levano uno strepito spaventato.

%%K t07-c2-14-2 p
Il \VUgras{gallo} \textsuperscript{(71)}, più bellicoso, alza la sua \VUgras{cresta}
\textsuperscript{(72)} rossa, scuote le penne della sua \VUgras{coda}
\textsuperscript{(73)} e lancia un sonoro «chicchirichì».

%%K t07-c2-15-1 p
15. --- Una \VUgras{chioccia} \textsuperscript{(74)} razzola sul \VUgras{letamaio}
\textsuperscript{(81)} con i suoi \VUgras{pulcini} \textsuperscript{(75)}. Il
\VUgras{porcellino} \textsuperscript{(78)} corre da sua madre, l'enorme
\VUgras{scrofa} \textsuperscript{(77)} che vediamo accanto al porcile.

%%K t07-c2-16-1 p
16. --- Nelle loro gabbie i \VUgras{conigli} \textsuperscript{(79)} mangiano avidamente
l'\VUgras{erba} \textsuperscript{(80)} tenera che porta loro la figlia del fattore.
Giannina vuole molto bene ai suoi conigli, e ogni giorno, dopo aver preso le
\VUgras{uova} \textsuperscript{(76)} fresche dal pollaio, viene a trovare i suoi
piccoli amici.
\VUsaut{6.0mm}
\VUfilet{22mm}
